% DRAFT, rewrite
\section{User interface design}

\subsection{A section's states}
Screens do not draw their own spinners or empty states for a part. A section's rows
are wrapped in \texttt{arriving(phase, hasRows:, title:)}, which draws exactly one of
these. Rows already on screen stay there whatever their part is doing: the status
replaces only rows that do not exist.

\begin{center}\small
\begin{tabularx}{\linewidth}{@{}l X X@{}}
  \toprule
  Phase & What is drawn & Accessibility \\
  \midrule
  \texttt{held} & The content. & Read as content. \\
  \texttt{late} & The content, under one line: ``Dati di 12 min fa, sto aggiornando''. The line goes when the answer arrives. & The line is read before the content. \\
  \texttt{waiting}, under 300 ms & The remembered height, empty. & Hidden. \\
  \texttt{waiting}, over 300 ms & The title and one soft waiting line where the first text line will be. & One element: ``Caricamento'' and the title. \\
  \texttt{waiting}, over 3 s & The same, plus the step in progress or ``120 di 340''. & The step is read once. \\
  \texttt{arrived} & The content, with the reveal if it is new. & Announced if the wait passed 3 s. \\
  \texttt{empty} & The section's empty-state message. & Read as content. \\
  \texttt{failed} & One line naming what is unavailable and ``Riprova''. & A button. \\
  \bottomrule
\end{tabularx}
\end{center}

A service that is paused after two timeouts is named once for the whole app, in
the status line: ``Carriera non risponde: mostro i dati salvati'', or the count
when there are several. Sections say nothing more than their \texttt{late} or
\texttt{failed} state.

\subsection{The reveal}
\begin{itemize}
  \item Text: a \texttt{TextRenderer} gives each group of five glyphs, about a word, opacity 0 to 1, blur 6 to 0 pt and a 2 pt rise, staggered so that the whole text ends within 0.6 s; only the first three lines animate piece by piece, the rest fade as one. The layout a renderer receives has glyphs and runs, not words.
  \item Rows and cards: opacity, a 6 pt blur and a 2 pt rise over 0.3 s, each starting 0.05 s after the one above and none later than 0.3 s, so the section ends within 0.6 s however long it is.
  \item Numbers: \texttt{.contentTransition(.numericText(value:))}.
  \item On a refresh over visible content, only the rows the refresh added or changed replay their reveal, and their mark comes in with it. Every other row stays as it is.
  \item The reveal plays only when rows arrive from the network in a section that was seen waiting. Rows that are there when the screen opens, a stored copy, or a day paged to are drawn at once.
  \item A tap or a scroll ends the reveal at once. Reduce Motion replaces it with a 0.2 s cross-fade and removes the waiting line's movement.
\end{itemize}

\subsection{New and changed items}
A row that a refresh added or changed carries a mark: a small dot in the accent
colour at its top trailing corner, read by VoiceOver as ``Nuovo'' after the row.
In the libretto a closed year carries the mark when one of its results does. The
marks of a section go away once the section has been on screen for at least a
second and the student leaves it or puts the app away; a section flicked past
keeps them. Appelli and Libretto in Carriera are the first sections with marks.

\subsection{Lists that grow}
A list that receives rows while the student has scrolled adds them below the
visible rows. A sorted list that would reorder waits for the scroll to stop, then
moves rows with the system's default list animation. Aule libere is the first such
list: while a campus is arriving the rooms shown keep their places, a line under
the header says how many rooms have answered, and ``nessuna aula libera'' is not
said until the whole campus has.

\subsection{Map and 3D}
Building colours on the map follow the order of arrival, which RelevanceOrder makes
the order of distance from what is on screen. A building's 3D shell fades in by
opacity when its model has loaded; no other effect is used, because shadows are
what costs frames on the iPhone 12 Pro.

\subsection{New pages in Impostazioni}
\begin{itemize}
  \item \textbf{Offline}: a switch for the pack, one row per item with ``pronto'' and its age or ``non ancora scaricato'', the space used, and ``Scarica ora''.
  \item \textbf{Dati e archiviazione}, which exists, gains a section ``Preparazione in anticipo'' rather than a page of its own: ``Prepara in anticipo'' on by default, ``Anche con la rete cellulare'' off by default and shown only while the first is on, and ``Richieste in anticipo oggi'' with today's count against the limit. The footer says what is prepared and when it is not.
  \item \textbf{Promemoria}, which exists: the switches for the kinds of news gain ``Lezioni spostate''; marks and sittings already had theirs. A new switch at the foot of the page, ``Vibrazione all'arrivo dei dati'', off by default, enables the haptic of \idref{R10}.
\end{itemize}

\subsection{Diagnostics}
The diagnostics report gains ``Caricamento'': per screen, the median and 90th
percentile of time to first content and to the last part; per kind of part, the
same for its duration and its share of failures; the share of speculative requests
used; and ``campione troppo piccolo'' under 20 openings.
